% !TeX root = proyecto.tex

%========================================================================================
% CAPÍTULO 7: DESARROLLO DEL TT2
%========================================================================================

\chapter{Desarrollo del TT2: Prototipos 3 y 4}
\label{cap:desarrolloTT2}

El presente capítulo vertebra el núcleo ingenieril y arquitectónico del Trabajo Terminal II. Tras haber documentado exhaustivamente las limitaciones operativas, algorítmicas y de persistencia inherentes a los prototipos preliminares, esta sección desglosa analíticamente la conceptualización, diseño, codificación y despliegue de los Prototipos 3 y 4. Estos artefactos representan la cristalización del protocolo de registro 2026-A135, materializando la transición definitiva de un sistema sintáctico-estático hacia una plataforma neuro-simbólica robusta, tolerante a fallos e impulsada por modelos fundacionales de Procesamiento de Lenguaje Natural (NLP).

\section{Diseño de Patrones y Arquitectura de Software}
\label{sec:tt2_patrones}

Para garantizar la mantenibilidad y escabilidad del sistema, se implementaron patrones de diseño consolidados a nivel empresarial, aislando responsabilidades y optimizando el consumo de recursos computacionales.

\subsection{Patrón Repository para la Persistencia}
\label{subsec:tt2_repository}

El acceso a la capa de datos relacional se desacopló por completo de la lógica de negocio mediante la implementación del Patrón \textit{Repository} (\texttt{NoticiasRepository} en \texttt{src/database/repository.py}). En los prototipos iniciales, las consultas de filtrado exigían cargar el archivo CSV íntegro en estructuras \textit{DataFrame} de Pandas, ocasionando asfixia de memoria RAM ($O(n)$) ante volúmenes masivos. 

Con el \textit{Repository}, el código delega el peso transaccional al ORM de SQLAlchemy. La búsqueda de subcadenas ya no depende de iteraciones forzadas; el repositorio explota los Índices Invertidos Generalizados (GIN) de PostgreSQL aplicados sobre un campo vectorial (\texttt{tsvector}). Esto garantiza que funciones analíticas como \texttt{buscar\_fts} resuelvan consultas heurísticas y de concurrencia masiva en rangos de latencia de milisegundos ($O(\log n)$), encapsulando de manera transparente la gestión de sesiones (\texttt{db.session}) y los fallbacks de tolerancia a fallos.

\subsection{Patrón Pipeline para el Flujo Analítico}
\label{subsec:tt2_pipeline}

El orquestador principal, \texttt{SimplifiedNewsAnalyzer}, materializa el Patrón \textit{Pipeline}. El flujo de ingesta de información no es monolítico, sino una secuencia determinista de estados discretos. Los datos extraídos transitan por estaciones de vectorización, reducción de dimensionalidad, similitud matricial e inferencia neuronal, permitiendo activar o desactivar módulos dinámicamente (\textit{Feature Flags}). Este diseño facilita una inyección de dependencias modular, donde cada paso (ej. \texttt{step\_9\_semantic\_detection}) recibe mutaciones del estado anterior sin violar el principio de responsabilidad única.

\section{Prototipo 3: Migración a la Arquitectura Semántica}
\label{sec:tt2_prototipo3}

El diseño del Prototipo 3 fue regido por el imperativo absoluto de erradicar la ceguera contextual del motor heurístico inicial. El componente lógico responsable de dotar de "comprensión" al sistema reside en el módulo \texttt{semantic\_detector.py}.

\subsection{El Motor de Detección BETO: Arquitectura de Inferencias}
\label{subsec:tt2_beto}

Se procedió a instanciar \textbf{BETO} (\textit{BERT para español}), un modelo transformador bidireccional entrenado sobre un vasto corpus hispanohablante (\textit{dccuchile/bert-base-spanish-wwm-cased}). El proceso interno transcurre en secuencias rigurosas:
Primero, el texto de la noticia es sometido a un algoritmo de fragmentación sub-palabra (\textit{WordPiece Tokenization}). Este mecanismo mitiga el problema crítico de las palabras fuera de vocabulario (\textit{Out-Of-Vocabulary}), descomponiéndolas en fonemas vectorizados reconocibles.

Una vez tokenizado y truncado a 512 secuencias posicionales, el arreglo ingresa en el paso hacia adelante (\textit{forward pass}) cruzando doce capas de mecanismos de autoatención (\textit{Self-Attention}). Estas capas obligan a cada \textit{token} a recalcular su valor en dependencia probabilística de su contexto izquierdo y derecho. La matriz resultante condensa el significado global en el token especial \texttt{[CLS]} como un hiper-vector de 768 dimensiones.

\subsection{Clasificación Zero-Shot y Temperature Scaling}

Para evadir la necesidad inicial de un entrenamiento supervisado masivo, el clasificador opera bajo el paradigma \textit{Zero-Shot}. Calcula la similitud del coseno entre la matriz de la noticia y tensores que describen categorías estáticas (Relevante / No Relevante). 

Sin embargo, las inferencias crudas entre espacios vectoriales tan complejos suelen generar puntajes muy estrechos (por ejemplo, 0.52 vs 0.48), provocando incertidumbre lógica. Para reducir drásticamente esta entropía, se implementó una función matemática de escalamiento termodinámico (\textit{Temperature Scaling}). Al inyectar un factor de escala térmico (multiplicador de 5) sobre el producto escalar normalizado antes de invocar a la exponencial euclidiana (\textit{Softmax}), se obliga a la distribución de probabilidad a separarse y polarizarse.

\begin{lstlisting}[language=Python, caption={Polarizaci\'on de probabilidades mediante Temperature Scaling}, label={lst:temp_scaling}]
# El multiplicador termico (5) amplifica la distancia matematica
# reduciendo la incertidumbre y forzando a la red a tomar una postura.
exp_rel = np.exp(sim_relevante * 5)
exp_norel = np.exp(sim_no_relevante * 5)

# Distribucion de la funcion Softmax polarizada
score = exp_rel / (exp_rel + exp_norel)
\end{lstlisting}

\section{Prototipo 4: Sistema Neuro-Simbólico y Agrupamiento Global}
\label{sec:tt2_prototipo4}

La creación del \textit{pipeline} neuro-simbólico responde a una necesidad imperiosa: amalgamar el razonamiento blando de las redes neuronales (conexionismo) con la rigidez innegociable de las reglas lógicas (simbolismo). 

\subsection{El Gran Filtro: BERTopic al inicio del Flujo}
\label{subsec:tt2_bertopic}

Históricamente en el TT1, la clasificación individual precedía al agrupamiento temático. En el Prototipo 4, la invocación del modelo BERTopic (\texttt{step\_10\_bertopic\_clustering}) fue relocalizada como paso primario de análisis. El objetivo de este rediseño es otorgarle a BETO un "Mapa Global" de la realidad periodística. 

En lugar de evaluar artículos en el vacío, BERTopic inyecta un identificador topológico (\texttt{topic\_id}) a cada nota. El proceso comprime la alta dimensionalidad con UMAP y detecta densidades con HDBSCAN. Esta cartografía permite al sistema reconocer \textit{outliers} (noticias que no se agrupan con ningún evento fáctico conocido). Cuando BETO identifica un \textit{outlier} (con \texttt{topic\_id = -1}), el sistema simbólico eleva automáticamente su umbral de tolerancia exigido (de 0.50 a 0.65) para considerarlo de Alta Prioridad. Así, el sistema exige un mayor rigor matemático a notas huérfanas, disminuyendo drásticamente los falsos positivos.

\subsection{Lógica Neuro-Simbólica: El Árbol Inquebrantable}
\label{subsec:tt2_logica_hibrida}

El pináculo del desarrollo reside en la función \texttt{step\_9\_semantic\_detection}. Para combatir las alucinaciones de la Inteligencia Artificial frente a contextos que replican semánticamente un crimen pero invierten los roles (e.g., menores agresores), el árbol de decisión interviene en tres estratos arquitectónicos.

\textbf{Nivel 1: El Bypass de Oro (Decisión Simbólica Absoluta).} \\
El modelo de Inteligencia Artificial carece de jurisdicción total cuando los algoritmos de recolección arrojan verdades analíticas absolutas. Si la recolección detecta expresiones explícitas de la taxonomía del dominio (\textit{Keywords de Oro}) asignando un puntaje de orfandad $\ge 0.80$, la lógica simbólica impone un \textit{override} incondicional. Se decreta como \textit{Alta Relevancia} sin importar la evaluación subyacente de BETO. Este diseño blinda al sistema de rechazar crímenes reales por problemas en la redacción periodística.

\textbf{Nivel 2: El Bozal de Penalización (Castigo Heurístico).} \\
Contrario al Bypass, el Bozal suprime inferencias equivocadas de BETO. Las noticias genéricas de violencia o aquellas donde el menor figura como el autor del homicidio activan la alerta temprana del sistema de recolección. Si el "score de feminicidio" o la relevancia heurística resultan nulas, el Bozal infiere un contexto nocivo. Multiplicando el valor probabilístico por un umbral fraccionario nulo, el sistema neutraliza la activación, asegurando que notas policiales irrelevantes jamás lleguen al Dashboard.

\textbf{Nivel 3: El Factor Alpha Dinámico.} \\
Para el grueso estadístico de notas ambiguas, se ejecuta la Fusión Matemática. Se genera un \textit{score\_final} balanceando el veredicto probabilístico de la red (\texttt{s\_score}) y la certidumbre extraída de variables duras (\texttt{h\_score}). El nivel de protagonismo que tiene BETO en esta ecuación es regido por la variable $\alpha$, cuya magnitud depende estrictamente de qué tan alejada está la red de la frontera de incertidumbre máxima (0.50). 

\begin{lstlisting}[language=Python, caption={Fusi\'on Neuro-Simb\'olica y c\'alculo del factor din\'amico Alpha}, label={lst:alpha_ponderation}]
# Distancia absoluta a la frontera de incertidumbre maxima (0.5)
distance = abs(s_score - 0.5)

if distance > 0.3:
    alpha = 0.7  # Distancia amplia: BETO tiene certeza, mayor peso
elif distance > 0.15:
    alpha = 0.5  # Distancia media: Ponderacion equitativa
else:
    alpha = 0.3  # BETO se encuentra indeciso: Delegar peso a la Heuristica

# Ecuacion de la recta con coeficientes dinamicos
score_final = (alpha * s_score) + ((1 - alpha) * h_score)

# Penalizacion direccional si la heuristica no concuerda
if s_score > h_score and distance > 0.2:
    score_final = max(h_score, score_final)
\end{lstlisting}

Mediante la sumatoria ponderada $\alpha$, el sistema confía ciegamente en BETO cuando la IA detecta correlaciones vectoriales contundentes (e.g. 0.98 o 0.12), pero se repliega a un comportamiento conservador fundamentado en la semántica dura de RegEx cuando la red titubea en el rango oscuro de 0.45 a 0.55.
